% Pais gap loop -- mathematics run sheet, references (loaded by \input from part 2). Each run appends its new bibitems here.
\begin{thebibliography}{99}\small
\bibitem{jackson} J.\,D. Jackson, \emph{Classical Electrodynamics}, 3rd ed.\ (Wiley, 1999), \S\S6.7, 9.2, 9.7--9.8.
\bibitem{laue} M. von Laue, ``Zur Dynamik der Relativit\"atstheorie'', Ann.\ Phys.\ \textbf{340}, 524 (1911).
\bibitem{moon} F.\,C. Moon, ``The virial theorem and scaling laws for superconducting magnet systems'' (Cornell, 1981); J.\ Appl.\ Phys.\ \textbf{53}, 9112 (1982).
\bibitem{hilal} M.\,A. Hilal, O. Arici, M. Cuban, ``Minimum structure magnetic energy storage solenoids with a toroidal insert'', IEEE Trans.\ Magn.\ \textbf{21}(2), 1044 (1985).
\bibitem{tsutsui} H. Tsutsui, S. Nomura, R. Shimada, ``Optimization of SMES coil by using virial theorem'', IEEE Trans.\ Appl.\ Supercond.
\bibitem{lbl} Lawrence Berkeley Laboratory report on force-compensated coils and the virial theorem, escholarship.org/content/qt0b08n8vh.
\bibitem{cr5} IEEE Council on Superconductivity, SMES review note CR5 (2008), snf.ieeecsc.org.
\bibitem{barth} C. Barth \emph{et al.}, ``Electro-mechanical properties of REBCO coated conductors from various industrial manufacturers at 77 K, self-field and 4.2 K, 19 T'', Supercond.\ Sci.\ Technol.\ \textbf{28}, 045011 (2015); arXiv:1502.06713.
\bibitem{jcsj} Teion Kogaku (J.\ Cryo.\ Soc.\ Jpn.) \textbf{46}(5), 2011: REBCO coil reversible to 1288 MPa / 0.63 \% hoop strain.
\bibitem{sandia} T. Lambert, ``Next generation composite materials for flywheel development'', DOE ESS peer review, Sandia (2012).
\bibitem{sae} ``Design, fabrication, and testing of 10 MJ composite flywheel energy storage rotors'', SAE 981282 (1998).
\bibitem{bai} Y. Bai \emph{et al.}, ``Carbon nanotube bundles with tensile strength over 80 GPa'', Nat.\ Nanotechnol.\ \textbf{13}, 589 (2018).
\bibitem{lee} C. Lee \emph{et al.}, Science \textbf{321}, 385 (2008).
\bibitem{ma} W. Ma \emph{et al.}, ``Monitoring a micromechanical process in macroscale carbon nanotube films and fibers'', Adv.\ Mater.\ (2009).
\bibitem{beyerlein} I.\,J. Beyerlein \emph{et al.}, ``Probabilistic strength of carbon nanotube yarns'' (Springer, 2010); Nanotechnology (2011).
\bibitem{gao} E. Gao, W. Lu, Z. Xu, ``Strength loss of carbon nanotube fibers explained in a three-level hierarchical model'' (2018).
\bibitem{daniels} H.\,E. Daniels, Proc.\ R.\ Soc.\ A \textbf{183}, 405 (1945).
\bibitem{rosen} B.\,W. Rosen, AIAA J.\ \textbf{2}, 1985 (1964).
\bibitem{pourrahimi} S. and N. Pourrahimi, US 6,534,718 and US 7,476,810, ``Reinforcement of superconducting coils by high-strength materials''.
\bibitem{cmut} ``Advances in capacitive micromachined ultrasonic transducers'', Micromachines \textbf{10}, 152 (2019).
\bibitem{ergun} A.\,S. Ergun \emph{et al.}, ``Broadband capacitive micromachined ultrasonic transducers ranging from 10 kHz to 60 MHz''.
\bibitem{coleman} S. Coleman, J.\,H. Van Vleck, Phys.\ Rev.\ \textbf{171}, 1370 (1968).
\bibitem{tajmar2} M. Tajmar \emph{et al.}, ``Thrust measurements of microwave-, superconducting- and laser-type EMDrives'' (TU Dresden).
\bibitem{tajmar1} M. Tajmar \emph{et al.}, ``SpaceDrive: thrust balance development'', IAC-18-C4.7-C3.5.5 (2018).
\bibitem{soni} J. Soni, S. Roy, ``Design and characterization of a nano-Newton resolution thrust stand'', Rev.\ Sci.\ Instrum.\ \textbf{84} (2013).
\bibitem{saulson} P.\,R. Saulson, ``Thermal noise in mechanical experiments'', Phys.\ Rev.\ D \textbf{42}, 2437 (1990).
\bibitem{schwinger} J. Schwinger, Phys.\ Rev.\ \textbf{82}, 664 (1951).
\bibitem{e320} E-320 collaboration talks: ELI (Sept 2025), ``first quantitative measurements of quantum radiation reaction''; LPHYS (2026); V. Yakimenko \emph{et al.}, Phys.\ Rev.\ Accel.\ Beams \textbf{22}, 101301 (2019).
\bibitem{e320p} ``Prospects for the production and detection of Breit--Wheeler tunneling positrons in Experiment 320'', arXiv:2506.04992 (2025).
\bibitem{gertsenshtein} M.\,E. Gertsenshtein, Sov.\ Phys.\ JETP \textbf{14}, 84 (1962).
\bibitem{bcs} J. Bardeen, L.\,N. Cooper, J.\,R. Schrieffer, Phys.\ Rev.\ \textbf{108}, 1175 (1957).
\bibitem{tinkham} M. Tinkham, \emph{Introduction to Superconductivity}, 2nd ed.\ (McGraw-Hill, 1996).
\bibitem{rowe} E. Rowe \emph{et al.}, ``Resonant enhancement of photo-induced superconductivity in K$_3$C$_{60}$'', Nat.\ Phys.\ \textbf{19}, 1821 (2023); arXiv:2301.08633.
\bibitem{shanghai} ``Numerical modeling of a high power terahertz source at Shanghai'', arXiv:1107.5619 (2011).
\bibitem{saser} C.\,L. Poyser \emph{et al.}, ``Terahertz sound lasers (sasers): recent developments and applications'', Proc.\ IOA; W. Maryam, PhD thesis, Nottingham (2013).
\bibitem{mankowsky} R. Mankowsky \emph{et al.}, Struct.\ Dyn.\ \textbf{4}, 044007 (2017); C.\,R. Hunt \emph{et al.}, Phys.\ Rev.\ B \textbf{94}, 224303 (2016).
\bibitem{mcculloch} M.\,E. McCulloch, ``Can the Podkletnov effect be explained by quantised inertia?'', arXiv:1108.3488 (2011).
\bibitem{dart} A.\,F. Cheng \emph{et al.}, Nature \textbf{616}, 457 (2023).
\bibitem{dessureault} Y.-S. Dessureault \emph{et al.}, ``Tensile strength and failure modes of continuous carbon nanotube yarns for structural composite application'', SAMPE 2020; Mater.\ Sci.\ Eng.\ A (2020).
\bibitem{liang23} ``Scalable high tensile modulus composite laminates using continuous carbon nanotube yarns for aerospace applications'', ACS Appl.\ Nano Mater.\ \textbf{6}(13), 11260 (2023), doi:10.1021/acsanm.3c01266.
\bibitem{wikiss} Manufacturer strength and density data as tabulated in ``Specific strength'', Wikipedia (accessed Nov 2026).
\bibitem{kellytyson} A. Kelly, W.\,R. Tyson, J.\ Mech.\ Phys.\ Solids \textbf{13}, 329 (1965).
\bibitem{ekin} J.\,W. Ekin, \emph{Experimental Techniques for Low-Temperature Measurements} (Oxford, 2006), thermal contraction tables.
\bibitem{scis22} ``Numerical research of potential methods to reduce screening-current-induced stress for no-insulation REBCO coils'', IEEE Trans.\ Appl.\ Supercond.\ (2022).
\bibitem{yan19} Y. Yan, C. Xin, Y. Tan, T. Qu, ``Screening current effect on the stress and strain distribution in REBCO high-field magnets'', arXiv:1909.07553.
\bibitem{talley} R.\,L. Talley, ``Twenty first century propulsion concept'', Phillips Laboratory PL-TR-91-3009 (1991).
\bibitem{tajmarbb} M. Tajmar, ``Biefeld--Brown effect: misinterpretation of corona wind phenomena'', AIAA J.\ \textbf{42}, 315 (2004).
\bibitem{fan13} K. Fan \emph{et al.}, ``Nonlinear terahertz metamaterials via field-enhanced carrier dynamics in GaAs'', Phys.\ Rev.\ Lett.\ \textbf{110}, 217404 (2013); arXiv:1208.1538.
\bibitem{fan12} K. Fan \emph{et al.}, ``Nonlinear metamaterials through enhanced impact ionization in GaAs at terahertz frequencies'', APS March Meeting 2012, Q17.11.
\bibitem{tifee} ``Terahertz radiation-induced sub-cycle field electron emission across a split-gap dipole antenna'', arXiv:1508.04737 (2015).
\bibitem{electromig} Report of THz-field-induced electromigration in sub-wavelength metallic gaps (single-cycle pulses to 400 kV/cm), core.ac.uk/download/43246148.
\bibitem{tate} J. Tate, B. Cabrera, S.\,B. Felch, J.\,T. Anderson, Phys.\ Rev.\ Lett.\ \textbf{62}, 845 (1989); Phys.\ Rev.\ B \textbf{42}, 7885 (1990).
\bibitem{tajmargm} M. Tajmar \emph{et al.}, ``Experimental detection of the gravitomagnetic London moment'', arXiv:gr-qc/0603033; M. Tajmar, C.\,J. de Matos, arXiv:gr-qc/0607086 (2006).
\bibitem{tajmargm2} ``Local photon and graviton mass and its consequences'', arXiv:gr-qc/0603032 (2006).
\bibitem{e320t} E-320 positron tracker preprint, arXiv:2604.23805 (April 2026).
\bibitem{toray23} Toray Industries, ``Toray develops TORAYCA T1200, the ultra high-strength carbon fiber'', press release, 29 Oct 2023 (1,160 ksi, 1.82 g/cm$^3$).
\bibitem{zhang13} R. Zhang, Y. Zhang, Q. Zhang, H. Xie, W. Qian, F. Wei, ``Growth of half-meter long carbon nanotubes based on Schulz--Flory distribution'', ACS Nano \textbf{7}, 6156 (2013).
\bibitem{zhu19} Z. Zhu \emph{et al.}, ``Rate-selected growth of ultrapure semiconducting carbon nanotube arrays'', Nat.\ Commun.\ \textbf{10}, 4467 (2019).
\bibitem{pugno06} N.\,M. Pugno, R.\,S. Ruoff, ``Nanoscale Weibull statistics'', J.\ Appl.\ Phys.\ \textbf{99}, 024301 (2006); arXiv:cond-mat/0504518.
\bibitem{pugno08} N.\,M. Pugno, F. Bosia, A. Carpinteri, ``Multiscale stochastic simulations for tensile testing of nanotube-based macroscopic cables'', Small \textbf{4}, 1044 (2008).
\bibitem{parkcaxis} M. Park \emph{et al.}, ``De-lamination characteristics of coated conductor for conduction cooled HTS coil'' (c-axis tensile strength 18--53 MPa).
\bibitem{delam} ``Delamination strengths of different types of REBCO-coated conductors and method for reducing radial thermal stresses of impregnated REBCO pancake coils'', IEEE Trans.\ Appl.\ Supercond.
\bibitem{kar19} S. Kar \emph{et al.}, ``Next-generation highly flexible round REBCO STAR wires with over 580 A\,mm$^{-2}$ at 4.2 K, 20 T'', Supercond.\ Sci.\ Technol.\ \textbf{32}, 10LT01 (2019).
\bibitem{aries} CERN Accelerating News, Bruker BHTS tape for H2020 ARIES: $J_E$ 600 A/mm$^2$, nearly 1000 A/mm$^2$ at 4.2 K, 20 T (2019).
\bibitem{selva24} ``40-meter-long REBCO tapes'' (University of Houston), 2024: short samples $J_e\sim4600$ A/mm$^2$ at 4.2 K, 20 T.
\bibitem{cooper70} R.\,K. Cooper, C.\,H. Woods, V.\,K. Neil, C.\,E. Taylor, ``Stability analysis of a levitated superconducting current ring stabilized by feedback and eddy currents'' (1970).
\bibitem{post} R.\,F. Post (LLNL), ambient-temperature passive magnetic bearings and Earnshaw's theorem, OSTI 632846; ``Earnshaw-stability of levitated rotors'', core.ac.uk 71164777.
\bibitem{anderson08} J.\,D. Anderson, J.\,K. Campbell, J.\,E. Ekelund, J. Ellis, J.\,F. Jordan, ``Anomalous orbital-energy changes observed during spacecraft flybys of Earth'', Phys.\ Rev.\ Lett.\ \textbf{100}, 091102 (2008).
\bibitem{iorio15} L. Iorio, ``Gravitational anomalies in the solar system?'', arXiv:1412.7673.
\bibitem{acedo18} L. Acedo, P. Piqueras, J.\,A. Mora\~no, ``A possible flyby anomaly for Juno at Jupiter'', Adv.\ Space Res.\ (2018); arXiv:1711.08893.
\bibitem{buhler} Exodus Propulsion Technologies (C. Buhler), charged-dielectric thrust claims: WO 2020/159603 A2 as cited in press reports; The Debrief and IFLScience coverage (claims, no independent test).
\bibitem{hathaway21} G. Hathaway, L.\,L. Williams, ``Null-result test for effect on weight from large electrostatic charge'', Physics \textbf{3}(2), 160 (2021); arXiv:2103.00700.
\bibitem{physt18} ``Carbon nanotubes share the load'', Physics Today \textbf{71}(7), 21 (2018), on Bai \emph{et al.}: about 43 GPa on the engineering cross-section.
\bibitem{jaroszynski25} J. Jaroszynski \emph{et al.}, ``Critical current, lengthwise fluctuations, and flux jumps in REBCO CC: a torque magnetometry study up to 45 T'', arXiv:2502.02706 (2025).
\bibitem{exodus891} US 11,511,891, ``System and method for generating forces using asymmetrical electrostatic pressure'' (Exodus Propulsion Technologies), citing Canning (2004).
\bibitem{tajmar15} M. Tajmar, G. Fiedler, ``Direct thrust measurements of an EM Drive and evaluation of possible side-effects'', AIAA 2015-4083 (TU Dresden).
\bibitem{apexld} ``Design and operation of APEX-LD: a compact levitated dipole for the confinement of electron--positron pair plasmas'', arXiv:2511.07191 (2025).
\bibitem{ldx07} ``First flight of the Levitated Dipole Experiment'', APS DPP 2007 (560 kg coil, 40 min first flight, digital feedback).
\bibitem{ldxfpa} D. Garnier for the LDX team, ``Report on the Levitated Dipole eXperiment'', Fusion Power Associates (2009): over 80 hours of flight.
\bibitem{genta} G. Genta \emph{et al.} (Politecnico di Torino), rotating and non-rotating damping in rotordynamics; eddy-current dampers and electrodynamic bearings (ISMB proceedings).
\bibitem{ismb10} ISMB-10 submission 78: forward whirl unstable above $\omega_{\lim}=\sqrt{k/m}\,(1+c_{nr}/c_r)$.
\bibitem{brandt93} E.\,H. Brandt, M. Indenbom, ``Type-II-superconductor strip with current in a perpendicular magnetic field'', Phys.\ Rev.\ B \textbf{48}, 12893 (1993).
\bibitem{toraym55j} Toray TORAYCA M55J fibre and unidirectional-composite data sheets (Toray CMA); M55J/EX-1515 laminate data (Toray TAC).
\bibitem{nickfilm} Strained bilayer nickelate films: arXiv:2608.01295; ``Structural modifications in strain-engineered bilayer nickelate thin films'', Nature \textbf{653}, 76 (2026); arXiv:2510.27613.
\bibitem{sun23} H. Sun \emph{et al.}, ``Signatures of superconductivity near 80 K in a nickelate under high pressure'', Nature \textbf{621}, 493 (2023).
\bibitem{nbh} Hydrogen-loaded epitaxial Nb films: about $-10$ GPa elastic stress below 6--8 nm (Pundt group, Göttingen and KIT).
\bibitem{post96} R.\,F. Post, D.\,D. Ryutov, ``The Inductrack: a simpler approach to magnetic levitation'' (LLNL); $L/D=2\pi vL/(\lambda R)$.
% ---- run 7
\bibitem{bai18si} Y. Bai \emph{et al.}, ``Carbon nanotube bundles with tensile strength over 80 GPa'', Nat. Nanotechnol. \textbf{13}, 589 (2018), Supplementary Information, Tables S2--S3 (areas $A_1$, $A_2$ and stresses $\sigma_1$, $\sigma_2$ per bundle).
\bibitem{lbc6} LBNL ``Little Big Coil'' LBC6, EUCAS 2025 submission (2025): Shanghai Superconductor EuBCO+BaHfO$_3$ tape, 4.0 mm $\times$ 45 $\mu$m (30 $\mu$m Hastelloy), $>1500$ A at 4.2 K, 30 T, $B\parallel ab$ for production-length tape; 42.6 T at first quench.
\bibitem{majkic20} G. Majkic \emph{et al.}, ``In-field critical current performance of 4.0 $\mu$m thick film REBCO conductor with Hf addition at 4.2 K and fields up to 31.2 T'', Supercond. Sci. Technol. (2020): 5812 A per 4 mm at 30 T, $B\parallel ab$.
\bibitem{exodusslides} Exodus Propulsion Technologies, company presentation (load-cell reading 4,460 $\mu$N at 0.065 $\mu$A, 3,500 V, $1.38\e{-5}$ Torr); status: okdiario, 22 June 2026, citing Popular Mechanics (May 2026).
\bibitem{stepan89} G. St\'ep\'an, \emph{Retarded Dynamical Systems: Stability and Characteristic Functions} (Longman, 1989): the D-subdivision method.
\bibitem{breda05} D. Breda, S. Maset, R. Vermiglio, ``Pseudospectral differencing methods for characteristic roots of delay differential equations'', SIAM J. Sci. Comput. \textbf{27}, 482 (2005).
\bibitem{freudenberg85} J.\,S. Freudenberg, D.\,P. Looze, ``Right half plane poles and zeros and design tradeoffs in feedback systems'', IEEE Trans. Autom. Control \textbf{30}, 555 (1985); H.\,W. Bode, \emph{Network Analysis and Feedback Amplifier Design} (1945).
\bibitem{stein03} G. Stein, ``Respect the unstable'', IEEE Control Syst. Mag. \textbf{23}(4), 12 (2003).
\bibitem{beams51} J.\,W. Beams, J.\,D. Ross, J.\,F. Dillon, ``Magnetically suspended vacuum-type ultracentrifuge'', Rev. Sci. Instrum. \textbf{22}, 77 (1951); J.\,W. Beams, double magnetic suspension with servo-stabilized vertical axis, Rev. Sci. Instrum. \textbf{34}, 1071 (1963).
\bibitem{humphreys09} D.\,A. Humphreys \emph{et al.}, ``Experimental vertical stability studies for ITER performance and design guidance'', Nucl. Fusion \textbf{49}, 115003 (2009); DIII-D vertical position control with vessel eddy currents (Nucl. Fusion, core.ac.uk 147937471).
\bibitem{esg} D.\,A. Hill, T. Letendre, H.\,A. Mills, ``Embedded, real-time DSP control of an electrostatically suspended gyroscope'', Proc. ACC 2004; US Patent 3,955,426 (mass-unbalance modulation removed from the suspension servo).
\bibitem{tran80} M.\,Q. Tran, A. Lee, R. Bollens, ``Passive feedback stabilization of one levitated coil by means of superconducting loops'', CRPP Lausanne report LRP 172 (1980) (record seen; full text not read).
\bibitem{powell69} J.\,R. Powell, G.\,T. Danby, null-flux suspension, US Patent 3,470,828 (1969); M. Pelovitz, W.\,J. Heim, J.\,G. Murray, ``Floating ring stabilization system'', IEEE Trans. Nucl. Sci. \textbf{18}(4), 91 (1971).
\bibitem{us5789837} US Patent 5,789,837, ``High-temperature superconducting magnetic bearing'' (superconducting shorted turns; flux creep in bulk bearings); US 5,481,146 (passive null-flux coil bearing).
\bibitem{nimsjoint} Y. Takeda \emph{et al.} (NIMS), slurry-processed REBCO joints below $10^{-12}\,\Omega$ at 4 K (EUCAS 2025); REBCO single-turn loop joint below $10^{-13}\,\Omega$ (IEEE Trans. Appl. Supercond., 2019); soldered joints: Mini-RT REBCO coil, 320 h decay at 25 K (Plasma Fusion Res. 9, 1405014, 2014).
\bibitem{jiang17} Z. Jiang \emph{et al.}, ``Dynamic resistance of a high-$T_c$ coated conductor wire in a perpendicular magnetic field at 77 K'', Supercond. Sci. Technol. \textbf{30}, 03LT01 (2017); M.\,D. Ainslie \emph{et al.}, Supercond. Sci. Technol. \textbf{31}, 074003 (2018).
\bibitem{st194} Dynamic resistance of YBCO-coated conductors in applied AC fields with DC transport currents (IEEE CSC slides ST194): $R_{\rm dyn}/L=4af(B_{ac}-B_{\rm th})/I_c$; perpendicular thresholds 2--3$\times$ the slab value.
\bibitem{felice} B.\,L. Militsyn \emph{et al.}, ``FELICE --- the free electron laser for intra-cavity experiments'', Nucl. Instrum. Methods A (2003); B. Redlich, A.\,F.\,G. van der Meer, FEL 2007 (10 $\mu$J outcoupled $\approx$ 400 $\mu$J intra-cavity); arXiv:2203.05348 (2--5\,\% hole outcoupling). Stations: molecular beam and FTICR cell (gas phase), Rayleigh range 55 and 82 mm (FELIX Laboratory wiki); 4-mirror 9 m resonator, 5 cm Rayleigh range at the experiment (FEL 2007, MOAAU04).
\bibitem{magsailrev} ``Propellantless space exploration'', arXiv:2510.21743 (2025), Sec.~V; R. Zubrin, D. Andrews, J. Spacecraft Rockets \textbf{28}, 197 (1991).
\bibitem{shaikh11} I.\,F. Shaikhislamov \emph{et al.}, mini-magnetosphere in the kinetic regime, arXiv:1110.4461 (2011): size equals the MHD standoff when $d_i<R_M$ and the St\"ormer radius otherwise.
\bibitem{khazanov04} G.\,V. Khazanov \emph{et al.}, ``Plasma sail concept fundamentals'', NASA/TP-2004-213143 (2004): fluid models invalid at M2P2 scales.
\bibitem{niquench22} Current distribution and induced edge currents in no-insulation REBCO coils when an adjacent coil quenches, J. Supercond. Nov. Magn. (2022), doi:10.1007/s10948-022-06344-z (one quenched NI pancake loses almost all its current; neighbours carry induced edge currents).
\bibitem{cernparallel} Protection of parallel-connected coils: the current of a quenched coil transfers to the coils in parallel through their inductive coupling; the quench stays local only if they stay below $I_c$ (CERN indico event 763185, contribution 3416442).
\bibitem{us10984934} US Patent 10,984,934, ``Fast inductive heaters for active quench protection of superconducting coil'' (in an NI coil a quenching region's current falls and nearby turns rise toward $I_c$ by inductive coupling).
\bibitem{lee26scs} J. Lee \emph{et al.}, ``Possible self-arresting screening current stress damage in a REBCO test coil at 42.6 T in 31 T background field'', IEEE Trans. Appl. Supercond. (2026) (EUCAS 2025); J. Bang \emph{et al.}, ``Elastic-plastic conductor damage evaluation at over 0.4\,\% strain using a high-stress REBCO coil'', Supercond. Sci. Technol. \textbf{37}, 095011 (2024) (Hastelloy-substrate tape yields at 0.9--1.1 GPa).
\bibitem{dyneema} DSM Dyneema SK75 technical brochure: tenacity 35.1 cN/dtex (3.4 GPa), modulus 113 GPa, elongation 3.5\,\%; SK staple 3.3--3.9 GPa.
\bibitem{flywheelcontain} A. J. Colozza, ``High energy flywheel containment evaluation'', NASA/CR-2000-210508 (2000); A. Buchroithner \emph{et al.}, Appl. Sci. \textbf{8}, 2622 (2018) (burst containment; housing dominates system mass); US Patent 6,889,577 (energy-absorbing housing; debris energy equals rotor energy).
\bibitem{vdlaan07} D. C. van der Laan, J. W. Ekin, C. C. Clickner, T. C. Stauffer, ``Delamination strength of YBCO coated conductors under transverse tensile stress'', Supercond. Sci. Technol. \textbf{20}, 765 (2007); H.-S. Shin, M. A. Diaz, MT28 (commercial tapes 22--38 MPa at 77 K); Fermilab/LBNL slides (2--100 MPa).
\bibitem{striated} R. Gyur\'aki \emph{et al.}, ``Interfilament resistance at 77 K in striated HTS coated conductors'', IEEE Trans. Appl. Supercond. \textbf{26} (2016) (10 filaments, 20 $\mu$m grooves); J. \v{S}ouc \emph{et al.}, ``AC losses of copper stabilized multifilament YBCO coated conductors'', arXiv:1212.0452 (stripes 0.04--0.07 mm); fully stabilized tapes with 12, 24, 46 filaments, $I_c$ 410 $\to$ 360, 310 A (Ohio State).
\bibitem{ldx} ``Status of the floating coil of the Levitated Dipole Experiment'', IEEE Trans. Appl. Supercond. \textbf{12}, 666 (2002) (6--8 h levitated as the coil warms from 5 to 10 K); A. Zhukovsky, D. T. Garnier, ``Operation of the LDX floating coil'', IEEE Trans. Appl. Supercond. \textbf{16}, 898 (2006); OSTI 781377.
\bibitem{rabenhorst} D. W. Rabenhorst, Johns Hopkins APL Technical Digest \textbf{1}(2) (flywheel energy densities: isotropic $\sim$12 Wh/kg in practice, wound 22--80 Wh/kg, $\sim$550 Wh/kg theoretical for a thin ring of the best materials).
\bibitem{tamuflywheel} Texas A\&M shaftless flywheel study (composite flywheel 42 kJ/kg with shaft; 5.6 Wh/kg for the whole system).
\bibitem{us5921505} US Patent 5,921,505, ``System and method for reducing mechanical disturbances from energy storage flywheels'' (counter-rotating pairs; residual torque sensed and nulled); US 6,232,671.
\bibitem{magnet300} S. Chen \emph{et al.}, ``Development of an 8-T conduction-cooled superconducting magnet with 300-mm warm bore for material processing application'', IEEE Trans. Appl. Supercond.; Magnex 9.4 T/310 mm specification TS1330H (2005).
\bibitem{dielectricsheets} PTFE: volume resistivity $>10^{18}\,\Omega$\,cm, $\varepsilon_r=2.1$ (Professional Plastics; Proto Labs data sheet 2026); polyimide film: $10^{15}\,\Omega$\,cm (CS Hyde) to $1.5\times10^{17}\,\Omega$\,cm, $\varepsilon_r=3.4$ (ERIKS).
\bibitem{elbe06} U. Lehnert \emph{et al.}, ``Laser gain and intra-cavity losses of the ELBE mid-IR FEL'', Proc. FEL 2006, TUPPH014 (round-trip loss 5.25--10\,\%, 6\,\% gain over 7\,\% loss); FELIX first lasing, Proc. EPAC 1992, p.~46 (single-pass gain of order 50\,\%).
\bibitem{dodge23} J. S. Dodge, L. Lopez, D. G. Sahota, ``Optical saturation produces spurious evidence for photoinduced superconductivity in K$_3$C$_{60}$'', Phys. Rev. Lett. \textbf{130}, 146002 (2023); reply arXiv:2303.10169; ``More proof that the evidence for photoinduced superconductivity in K$_3$C$_{60}$ is distorted'', APS March Meeting 2025.
\bibitem{cavitysc} H. Zhang \emph{et al.}, ``Cavity enhanced superconductivity'', arXiv:2606.19171 (2026) (NbSe$_2$, 3.02 $\to$ 3.41 K); vacuum-dressed NbN in a high-$Q$ THz cavity, arXiv:2601.08191; cavity-altered superconductivity ($\kappa$-ET next to hBN), arXiv:2505.17378; A. Thomas \emph{et al.}, ``Exploring superconductivity under strong coupling with the vacuum electromagnetic field''; QEDFT for MgB$_2$, arXiv:2404.08122.
\bibitem{nogo26} ``No-go theorem and routes towards cavity-enhanced superconductivity'', arXiv:2608.14784 (2026); Casimir-cavity $T_c$ test, no shift above 12 $\mu$K.
% ---- run 9
\bibitem{virialgas} The virial theorem for a gas in a vessel, $2\langle E_{\rm kin}\rangle=3pV$: e.g.\ ``What is osmosis?'', arXiv:physics/0305011, App.~A; with $P=(\gamma-1)u$, $3(\gamma-1)E_i+\Omega=0$ (coefficient 2 for $\gamma=5/3$): ``Axion astrophysics'', arXiv:2411.02492, \S2.6.1.
\bibitem{rotliner} Rotatable flywheel containment: US 6,889,577 (liner released above a torque threshold, then riding on a fluid film); US 11,143,277 (free-spinning liner absorbs the burst rotor's angular momentum); UT-CEM full-scale composite containment for the ATTB (SAE 2004-01-0005; burst at 39,500 rpm, housing torque $\le$ 63,000 ft-lb against 400,000 predicted, spin-down 120 s; core.ac.uk 322358635).
\bibitem{whipple} E.\,L. Christiansen, modified Cour-Palais Whipple-shield equation for $V_n\ge7$ km/s (NASA JSC memo, NTRS 19920010826); ballistic, shatter and melt regimes below 3, 3--7 and above 7 km/s (Acta Astronaut., 2024, Padova); stuffed Whipple shields (zenodo 1258555).
\bibitem{blastvessel} M.\,A. Polcyn, E.\,D. Esparza, M.\,G. Whitney, ``Design of the M-9 firing facility containment vessel for Los Alamos'' (10 kg TNT, weight limit 75 t; DTIC ADA525803) and its proof test (DTIC ADA499952); LLNL, ``Composite vessels for containment of extreme blast loadings'' (2 m vessel for 35 kg TNT; OSTI 15014483); Sandia EDS vessels (11 kg TNT, 1092 shots).
\bibitem{ipacs} C.\,M. Roithmayr \emph{et al.}, ``Integrated power and attitude control for a spacecraft with flywheels and control moment gyroscopes'', NASA/TP-2003-212178; C.\,M. Roithmayr, ``International Space Station attitude motion associated with flywheel energy storage'' (abrupt stop of one rotor of a counter-rotating pair), NTRS 20040086750; US 6,113,033.
\bibitem{senatore} C. Senatore (Univ.\ Geneva), HiTAT workshop slides, CERN indico 1220254: $I_c$/width 1094--1218 A/cm at 4.2 K, 19 T and 341--460 A/cm at 20 K, 19 T ($B\perp$); $\alpha=0.75$ above 10 T; $I_{c\parallel}/I_{c\perp}=3.5$--3.8; $T^*=21$--23 K below 40 K.
\bibitem{selvafnal} V. Selvamanickam, ``Conductor and cable challenges and opportunities'', FNAL (2023): 50 m A-MOCVD tape at 3.5$\times$ the $I_c$ of typical commercial tape at 4.2 K, 20 T; STAR wires $>2.5$ kA at 4.2 K, 30 T.
\bibitem{htsdb} S.\,C. Wimbush, N.\,M. Strickland, ``A public database of high-temperature superconductor critical current data'', IEEE Trans.\ Appl.\ Supercond.\ \textbf{27}, 8000105 (2017); figshare collection 2861821 ($I_c(T,B,\theta)$, 15--90 K, up to 8 T).
\bibitem{lpf} M. Armano \emph{et al.} (LISA Pathfinder), ``Charge-induced force noise on free-falling test masses'', Phys.\ Rev.\ Lett.\ \textbf{118}, 171101 (2017); ``Precision charge control for isolated free-falling test masses'', Phys.\ Rev.\ D \textbf{98}, 062001 (2018); charge modulation by phased UV pulses, Phys.\ Rev.\ D \textbf{106}, L101101 (2022); D. Dal Bosco, PhD thesis (Trento): torsion-pendulum tests of the LISA UV-LED charge management.
\bibitem{fhifel} W. Sch\"ollkopf \emph{et al.}, ``First lasing of the IR FEL at the Fritz Haber Institute'', FEL 2012, MOOB01 (computed small-signal gain $>200$\,\% for 1 ps bunches at 7--30 $\mu$m; 1\,\% loss per mirror).
\bibitem{elbeu100} ELBE far-IR FEL (U100): 10\,\% measured loss per pass at 30 $\mu$m with a 2 mm hole, FEL 2007 MOPPH036; design gain $>25$\,\% across the range, losses $\le6$\,\% without outcoupling, FEL 2005 TUPP031.
\bibitem{k3c60films} K$_3$C$_{60}$ on Si(111) by MBE, 1--140 ML at about 2 ML/h, superconducting down to $\ge25$ ML (TU Delft, Phys.\ Rev.\ B, 2003; arXiv:cond-mat/0206221); MBE growth of A$_x$C$_{60}$ films, AAPPS Bull.\ (2021), arXiv:2202.05968; Rb$_3$C$_{60}$ films 3--9 ML, $T_c$ up to 23 K, arXiv:2102.07914.
\bibitem{wang23} E. Wang \emph{et al.}, ``Superconducting nonlinear transport in optically driven high-temperature K$_3$C$_{60}$'' (2023): thin films on photoconductive switches, granular superconducting-like response.
\bibitem{devecchi25} G. de Vecchi \emph{et al.}, ``Upper limit on magnetic field expulsion in optically driven K$_3$C$_{60}$'', Phys.\ Rev.\ Research \textbf{7}, 043270 (2025); arXiv:2502.20276 (about 50 nT resolution, $\chi_v>-0.1$).
\bibitem{dfgmir} Mid-IR difference-frequency sources: C. Gaida \emph{et al.}, Light Sci.\ Appl.\ \textbf{7}, 94 (2018) (watt-scale intrapulse DFG); 6.7 W at 2.9 $\mu$m, Opt.\ Lett.\ \textbf{45}, 1248 (2020); none at watt level near 30 $\mu$m.
% ---- run 10
\bibitem{iseult} N. Boulant, L. Quettier \emph{et al.}, ``Commissioning of the Iseult CEA 11.7 T whole-body MRI'', MAGMA \textbf{36}, 175 (2023) (NbTi, 1.8 K, 338 MJ, 132 t; spec $<0.5$ ppm peak to peak on a 22 cm sphere, 0.05 ppm/h); ISMRM 2022 abstract 0382 (measured 0.88 ppm peak to peak, RMS 0.15 ppm, passive and cryo shims).
\bibitem{birb} T. Copp, J.\,G. Tuttle \emph{et al.}, ``The area-density-dependence of Ball IR Black's low-temperature emissivity'', Cryogenics \textbf{64}, 240 (2014) ($\varepsilon\ge0.95$ from 30 to 300 K); J. Tuttle \emph{et al.}, NASA NTRS 20180000042 (most black paints fall steeply below about 100 K).
\bibitem{acktar} Acktar Fractal Black, total hemispherical emissivity 5--300 K, TFAWS 2023 (TFAWS23-CR-1): 0.8--0.9 at 40--300 K, falling significantly below about 50 K; 30 $\mu$m effective thickness.
\bibitem{rebconmr} J. Bang \emph{et al.}, ``Harmonic errors of a 9.4 T all-REBCO NMR magnet affected by screening current and geometric inconsistency of coated conductors'', Sci.\ Rep.\ (2024), PMC11333634 (over 100 ppm between design and measurement; sub-ppm after passive and active shimming, 2018).
\bibitem{scfsim} S. Shi, R. Liang, ``Numerical analysis of REBCO high-temperature superconducting coils based on screening effect'', J.\ Supercond.\ Nov.\ Magn.\ \textbf{35}(12) (2022) (screening-current field up to 0.8 T, 7.8\,\%); current overshoot and field shaking remedies, J. Korea Inst.\ Appl.\ Supercond.\ Cryog.\ (2018).
\bibitem{schonberg} W.\,P. Schonberg, R.\,A. Taylor, J.\,R. Horn, ``An analysis of penetration and ricochet phenomena in oblique hypervelocity impact'', NASA TM (MSFC, 1988); ``Oblique hypervelocity impact response of dual-sheet structures'' (NTRS 19890011875; largest ricochet fragment up to 40\,\% of the projectile diameter at 24--36\,\% of its speed); W.\,P. Schonberg, A.\,R. Ebrahim, Int.\ J.\ Impact Eng.\ \textbf{23}, 823 (1999).
\bibitem{swrilgg} Southwest Research Institute, ``Large two-stage light-gas gun facility'' (38 mm launch tube, 22 m, 7 km/s top speed, 0.5 kg at 3 km/s, 1.4 kg at 2 km/s; fragment and flyer-plate launches); JHU HyFIRE (1--7 km/s); Illinois HIT lab ($>7$ km/s).
\bibitem{wilsonmag} M.\,N. Wilson, \emph{Superconducting Magnets} (Oxford, 1983), ch.~7 (adiabatic stability of a slab, $\beta<3$).
\bibitem{kittel} C. Kittel, \emph{Introduction to Solid State Physics} (Wiley), ch.~5--6 (Debye $T^3$ and electronic $\gamma T$ heat capacity; Ni $\gamma\approx7.0$ mJ mol$^{-1}$K$^{-2}$).
\bibitem{tfmc} Z.\,S. Hartwig \emph{et al.}, ``The SPARC Toroidal Field Model Coil Program'', IEEE Trans.\ Appl.\ Supercond.\ \textbf{34}, 0600316 (2024), arXiv:2308.12301 (20.1 T at 40.5 kA, 815 kN/m, 10,058 kg, 270 km of REBCO, 600 W at 20 K, local damage after an open-circuit quench at 31.5 kA); B. LaBombard, IEEE CSC SNF 57 (2024) (110 MJ, not self-protecting).
\bibitem{hfgwcav} L. Fischer \emph{et al.}, ``First characterisation of the MAGO cavity'', Class.\ Quantum Grav.\ \textbf{42}, 115015 (2025); Y. Kim \emph{et al.}, ``Search for high-frequency gravitational waves via re-analysis of cavity axion data'', arXiv:2511.17817 ($h_0\approx3.9\e{-21}$ near 5.3 GHz); cavity reach $h\sim10^{-21}$ at 250--330 MHz and $10^{-20}$ at 2.5--3.4 GHz (Valero \emph{et al.}, 2024).
\bibitem{brack22} T. Brack \emph{et al.}, ``Measurement and theory of gravitational coupling between resonating beams'', arXiv:2201.03420 (2022) (42 Hz; detector amplitude up to $10^{-11}$ m; $2.5\e{-20}$ J/s; $d^{-4}$; 3\,\% agreement).
\bibitem{tajmar24} M. Tajmar, M. K\"o{\ss}ling, O. Neunzig, ``In-depth experimental search for a coupling between gravity and electromagnetism with steady fields'', Sci.\ Rep.\ \textbf{14} (2024), doi:10.1038/s41598-024-70286-w; arXiv:2402.15640.
\bibitem{zhang25pat} X. Zhang, Z. Xiao, X. Li, ``Method and device for driving movement of target object by gravitational field generated by varying electromagnetic field'', US Patent 12,417,870 B1 (16 Sep 2025).
\bibitem{patents25} CPC B64G1/417 filings: US 2024/0063731, US 2024/0375795, US 2025/0132082, US 2025/0286480, US 2025/0296708, US 12,442,361 (ipiq.io listing); US 2025/0055389 (``Voorhees Anti-Gravity Engine II'').
\bibitem{oldpatents} US 3,095,167 (electrostatic charge against Earth's charge); US 3,610,971 (motional electric field generator); US 3,120,363 (Electronatom lifter); US 3,071,705 (electrostatic fluid propulsion); US 2003/0209637 A1 (rotating electrostatic field; abandoned).
% ---- run 11
\bibitem{cyclotrons} J. Pradhan \emph{et al.}, ``Study of magnetic field imperfections of Kolkata superconducting cyclotron'', Cyclotrons 2010, MOPCP078 (a 2--3 G first harmonic disturbs the beam); LNS Catania first-harmonic correction (EPAC 1994, 2301); NSCL K500 (Cyclotrons 1998, G-18); J.\,W. Kim \emph{et al.}, RIKEN SRC trim coils (PAC 1997, 6P013: errors below $\pm$15 G); D. Zhang \emph{et al.}, CYCIAE-230 trim rods (Cyclotrons 2016, MOP06: 0.24 G per mm).
\bibitem{ubsc} R. Kumagai \emph{et al.}, ``Fabrication and experiments on a 1/2-scale demonstration NI-REBCO coil system of a skeleton cyclotron'', IEEE Trans.\ Appl.\ Supercond.\ \textbf{34}(5) (2024); S. Noguchi; H. Ueda \emph{et al.}, ASC 2024 (CERN indico 1431972): quench at 250 A of 540 A, coil damage, screening-current and contraction stress; SCIF and overshoot, IEEE Trans.\ Appl.\ Supercond.\ \textbf{32}, doi:10.1109/TASC.2022.3171195 (2022).
\bibitem{gouging} R.\,D.\,M. Tachau, ``An investigation of gouge initiation in high-velocity sliding contact'' (Sandia, 1991): rails above 1.5 km/s, gun barrels above 4 km/s; C.\,H. Yew, T.\,G. Trucano, R.\,D.\,M. Tachau, Int.\ J.\ Impact Eng.\ \textbf{17}, 825 (1995); J.\,D. Cinnamon, \emph{Hypervelocity Gouging Impacts} (AIAA, 2009); P.\,M. Bassett \emph{et al.}, Wear \textbf{310}, 41 (2014).
\bibitem{fsp} Fragment-simulating projectiles in a two-stage light-gas gun with separating sabots (OSTI 2004149): mean apparent pitch and yaw 4.7--5.5$^\circ$ end-on, 32--43$^\circ$ side-on.
\bibitem{avalanche} Suppression of flux avalanches by metal layers: NbN films half-coated with Cu, arXiv:1211.3787; thermomagnetic instability with an adjacent metal layer, arXiv:1304.5405; J. Jaroszynski, CCA 2025 (centre-slit tapes stable, edge-slit tapes fluctuate).
\bibitem{brack23} T. Brack \emph{et al.}, ``Dynamic measurement of gravitational coupling between resonating beams in the hertz regime'', Nat.\ Phys.\ \textbf{18}, 952 (2022); ``Dynamic gravitational excitation of structural resonances in the hertz regime using two rotating bars'', Commun.\ Phys.\ \textbf{6}, 270 (2023) (detector $Q\approx10^4$, 42 Hz, $\sim$0.1\,\% uncertainty).
\end{thebibliography}
